\chapter{Modified gravity and dark-sector simulations}\label{ch:mg-simulations}
\chapterstatus{written}{2026-11-20}

SimulationCode separates the \emph{cosmology} (the background, \param{cosmology_model}) from the \emph{gravity model}
(the growth of perturbations and the force, \param{gravity_model}), so that one can be changed while the other is
fixed. This chapter describes the models, how the fifth forces of modified gravity are computed with and without
screening, and how dark-sector interactions enter the particle dynamics. The linear theory of the same models is in
\cref{ch:extended-models}; the quasi-static limits (\cls{QuasiStaticGravity}) are shared.

\begin{modulebox}{SimulationCode models}
\file{src/Cosmology.h} (maps \param{cosmology_model} to the background registry), \file{src/GravityModel.h} (base:
$G_{\rm eff}(k,a)/G$, growth factors, force), \file{GravityModel_Screened.h} (fifth force: linear, screening,
exact), \file{GravityModel_GR.h}, \file{_fofr.h}, \file{_DGP.h}, \file{_CubicGalileon.h}, \file{_Horndeski.h},
\file{_JBD.h}, \file{_symmetron.h}, \file{_dilaton.h}, \file{_mbeta.h}, \file{_mu.h}, \file{_CoupledQuintessence.h},
\file{GravityModelSettings.h}, \file{SimulationIDE.h}, \file{SpeciesTimeStepping.h}; the frontend
\file{SimulationCode/scripts/horndeski_frontend.py}; the multigrid solvers of \cref{ch:multigrid}.\par
Documentation: \file{SimulationCode/MODELS.md}, \file{SPECIES.md}.
\end{modulebox}

\section{Cosmologies and gravity models}\label{sec:mgsim-models}

\begin{table}[htbp]
\centering\small
\caption{Cosmologies (\param{cosmology_model}). All have CDM, baryons, photons, massless and massive neutrinos,
curvature, $\Lambda$ and the primordial parameters; the backgrounds are those of \cref{ch:background}.}
\label{tab:mgsim-cosmologies}
\begin{tabularx}{\textwidth}{@{}lX@{}}
\toprule
\code{LCDM}, \code{w0waCDM} & $\Lambda$CDM; CPL dark energy \\
\code{JBD} & Jordan--Brans--Dicke ($\omega_{\rm BD}$, $G_{\rm eff}/G$ today) \\
\code{DGP} & self-accelerating DGP (with $\Lambda$CDM: the normal branch is a gravity model) \\
\code{Quintessence}, \code{EDE} & scalar-field dark energy (exponential, inverse power law, PNGB, hilltop); early dark energy \\
\code{Tabulated} & $H(z)$ from a file (\code{H}, a CLASS/hi\_class \code{background.dat}, or $w(a)$) \\
\code{DarkScattering}, \code{CovariantDarkScattering} & $w$CDM with dark matter--dark energy momentum exchange \\
\code{CoupledQuintessence} & field coupled to CDM, $m\propto e^{-\beta\phi}$ \\
\code{FadingDarkSector} & the same with $\lambda=c$, $\beta=c'$ (\param{cosmology_fds_c}, \param{cosmology_fds_cprime}) \cite{ObiedVafaWu2026} \\
\code{IDE} & energy exchange $Q=\xi H\rho_{\rm DE}$ \\
\code{Horndeski} & reduced Horndeski theory from the frontend table \\
\code{CubicGalileon} & the tracker (no free parameter) \\
\bottomrule
\end{tabularx}
\end{table}

\paragraph*{Cosmology parameters.} Beyond the $\Lambda$CDM parameters each cosmology reads its own:
$w(a)=w_0+w_a(1-a)$ from \param{cosmology_w0}, \param{cosmology_wa}; for \code{JBD} $\omega_{\rm BD}$
(\param{cosmology_JBD_wBD}), $G_{\rm eff}/G$ today (\param{cosmology_JBD_GeffG_today}) and which $G$ defines the density
parameters (\param{cosmology_JBD_density_parameter_definition}: \code{hi-class} by default, or \code{Gbare},
\code{Gnewton}, \code{Gtoday}); for \code{DGP} $\Omega_{rc}=1/(4r_c^2H_0^2)$ (\param{cosmology_dgp_OmegaRC}). Quintessence has
$V=V_0u(\phi)$ with $V_0$ fixed by $\Omega_{\rm DE}$ today and $u$ chosen by \param{cosmology_quintessence_potential}
(\code{exponential}: \param{cosmology_quintessence_lambda}; \code{inverse_power_law}: \param{cosmology_quintessence_alpha};
\code{pngb}: \param{cosmology_quintessence_f}; \code{hilltop}: \param{cosmology_quintessence_mu},
\param{cosmology_quintessence_p}), starting from \param{cosmology_quintessence_phi_ini} and
\param{cosmology_quintessence_dphidlna_ini} deep in the radiation era; early dark energy has
$V=m^2f^2(1-\cos\phi/f)^n$ with the peak fraction \param{cosmology_ede_fEDE} at \param{cosmology_ede_zc}, the initial
angle \param{cosmology_ede_theta_ini} and \param{cosmology_ede_n}. \code{Tabulated} reads \param{cosmology_tabulated_filename}
in the format \param{cosmology_tabulated_format} (\code{H}: $z$, $H/H_0$; \code{CLASS}: a background file; \code{w}: $a$,
$w(a)$). Dark scattering has $w$ (\param{cosmology_ds_w}), the dark-energy sound speed \param{cosmology_ds_cs2} (for initial
conditions from the solver) and $\xi$ (\param{cosmology_ds_xi}) or, covariantised, $\alpha$ (\param{cosmology_cds_alpha});
coupled quintessence $\lambda$ and $\beta$ (\param{cosmology_cq_lambda}, \param{cosmology_cq_beta}); interacting dark
energy $w$, $\xi$ and the sound speed (\param{cosmology_ide_w}, \param{cosmology_ide_xi}, \param{cosmology_ide_cs2}; the
model needs $\xi/[3(1+w)]\le0$); and \code{Horndeski} the frontend's table (\param{cosmology_horndeski_filename}).

\begin{table}[htbp]
\centering\small
\caption{Gravity models (\param{gravity_model}).}
\label{tab:mgsim-gravity}
\begin{tabularx}{\textwidth}{@{}lXl@{}}
\toprule
Model & $G_{\rm eff}/G$ and nonlinear physics & Screening \\
\midrule
\code{GR} & 1 & --- \\
\code{f(R)} & Hu--Sawicky \cite{HuSawicki2007}, $\frac{1}{1+f_R}\left[1+\frac13\frac{k^2}{k^2+a^2m^2}\right]$ & chameleon \cite{KhouryWeltman2004} \\
\code{DGP} & $1+1/(3\beta)$, brane-bending mode \cite{KoyamaMaartens2006} & Vainshtein \cite{Vainshtein1972} \\
\code{CubicGalileon} & DGP form with $r_c(a)$, $\beta(a)$ on the tracker \cite{Barreira2013} & Vainshtein \\
\code{Horndeski} & $\frac{1}{2G_4}+\frac{1}{3\beta}$ from the table \cite{Kimura2012,Wright2023} & Vainshtein ($G_{3X}\neq0$) \\
\code{JBD} & $G\phi_*/\phi$ & --- \\
\code{Symmetron} & $1+\frac{2\beta^2(a)k^2}{k^2+a^2m^2}$ after $a_{\rm ssb}$ \cite{HinterbichlerKhoury2010,Brax2012} & potential \\
\code{Dilaton} & as symmetron, $m=m_0a^{-r}$ \cite{Brax2012} & potential \\
\code{mbeta} & as symmetron, power laws $\beta_0a^{n}$, $m_0a^{m}$ & potential \\
\code{mu} & $\mu(a,k)$: DE form, Bertschinger--Zukin \cite{BertschingerZukin2008} or a table & none \\
\code{CoupledQuintessence} & CDM feels $G(1+2\beta^2)$ from CDM, friction $\beta\phi'$ & --- \\
\code{FadingDarkSector} & as \code{CoupledQuintessence}, $\beta=c'$ & --- \\
\bottomrule
\end{tabularx}
\end{table}

The base \cls{GravityModel} solves the LPT growth equations to third order for a general $G_{\rm eff}(k,a)/G$
(\cref{ch:growth-transfer}) and computes the force $\nabla^2\phi=\frac32\Om_ma(\delta+\delta_{\rm FF})$, where the
\emph{fifth-force density} $\delta_{\rm FF}$ carries all modifications; a new model normally only provides
$G_{\rm eff}$, its quasi-static theory and optionally an exact solver. The $\mu$ model with the Bertschinger--Zukin form
$\mu=(1+\beta_1\lambda_1^2k^2a^s)/(1+\lambda_1^2k^2a^s)$ or a table $\mu(a)$ or $\mu(a,k)$ (e.g.\ from hi\_class)
modifies only the potential that moves the particles (not the lensing potential); a $k$-dependent $\mu$ gives
scale-dependent growth.

\paragraph*{Gravity-model parameters.} $f(R)$: $f_{R0}$ and $n$ (\param{gravity_model_fofr_fofr0},
\param{gravity_model_fofr_nfofr}); DGP: $r_cH_0/c$ (\param{gravity_model_dgp_rcH0overc}); symmetron: $a_{\rm ssb}$, $\beta$,
the vacuum range $L$ in $\Mpch$ and the powers of $V=-(\phi/\phi_*)^N/N+(\phi/\phi_*)^M/M$
(\param{gravity_model_symmetron_assb}, \param{gravity_model_symmetron_beta}, \param{gravity_model_symmetron_L_mpch},
\param{gravity_model_symmetron_N}, \param{gravity_model_symmetron_M}); dilaton: $A_2$, $\beta_0$, $r$ and $\xi=H_0/m_0$ (the
range today is $2998\,\xi\Mpch$; \param{gravity_model_dilaton_A2}, \param{gravity_model_dilaton_beta0},
\param{gravity_model_dilaton_r}, \param{gravity_model_dilaton_xi}); $(m,\beta)$: $(\beta_0,n,m_0,m)$ of $\beta=\beta_0a^n$,
$m=m_0H_0a^m$ (\param{gravity_model_mbeta_params}); $\mu$: the form \param{gravity_model_mu_parametrization} (\code{DE}:
$\mu=1+\mu_0\Omega_{\rm DE}(a)/\Omega_{\rm DE}(1)$ with \param{gravity_model_mu_mu0}; \code{BZ}: $(\beta_1,\lambda_1,s)$ in
\param{gravity_model_mu_BZ_params}; \code{table}: \param{gravity_model_mu_filename}, rows $a$, $\mu$ or $a$, $k$, $\mu$ on a
regular grid); Horndeski: the frontend's table (\param{gravity_model_horndeski_filename}, by default that of the
cosmology).

\section{Fifth forces and screening}\label{sec:mgsim-screening}

Screened theories suppress the fifth force in dense regions, which is essential for their nonlinear predictions. The
fifth-force density is computed in one of four ways:
\begin{enumerate}
\item \textbf{Linear}: $\delta_{\rm FF}(\vk)=\left[G_{\rm eff}(a,k)/G-1\right]\delta(\vk)$. Fast, exact on large scales,
      too much force in haloes.
\item \textbf{Approximate screening} \cite{WintherFerreira2014,Brando2023} (\param{gravity_model_screening}): the
      linear fifth force multiplied, in real space, by a screening factor $0\le f\le1$ of the local field:
      \begin{itemize}
      \item \emph{potential (chameleon-like)}: $f=\min\!\left(1,\ \epsilon\,\Phi_{\rm crit}(a)/|\Phi_N(\vx)|\right)$, with
            the Newtonian potential $\Phi_N$ and the thin-shell threshold $\Phi_{\rm crit}=\frac32|f_R(a)|$ for $f(R)$,
            $\phi(a)/(2\beta\Mpl)$ for the symmetron and $(m,\beta)$, $1/(2A_2)$ for the dilaton (\cref{sec:ext-quasistatic});
      \item \emph{density (Vainshtein)}: for a spherical top hat the field equation of the brane-bending mode gives
            $f=2(\sqrt{1+e}-1)/e$ with $e\propto(r_cH_0/c)^2\Om_m\,\delta/\beta^2$, evaluated on the density smoothed with a
            filter on a scale $R$ in units of the box (\param{gravity_model_dgp_smoothing_filter},
            \param{gravity_model_dgp_smoothing_scale_over_boxsize}; for the cubic Galileon and Horndeski
            \param{gravity_model_cubic_galileon_smoothing_filter}, \param{gravity_model_cubic_galileon_smoothing_scale_over_boxsize},
            \param{gravity_model_horndeski_smoothing_filter}, \param{gravity_model_horndeski_smoothing_scale_over_boxsize});
            for DGP the
            form with $\Om_m$ today and $1+\delta$ is used, which reproduces the exact solution best (0.5\% rms in $P(k)$
            at $z=0$);
      \end{itemize}
      with the efficiency $\epsilon$ (\param{gravity_model_screening_efficiency}) as a tuning parameter.
\item \textbf{Screening plus the linear large scales}:
      $\delta_{\rm FF}\to\delta_{\rm FF}^{\rm lin}\,F+\delta_{\rm FF}^{\rm screened}(1-F)$ with the low-pass filter
      $F=e^{-k^2/2k_*^2}$ (\param{gravity_model_screening_linear_scale_hmpc}), so that the linear evolution on large
      scales is exact (switch: \param{gravity_model_screening_enforce_largescale_linear}). Usually the best choice.
\item \textbf{Exact}: the nonlinear field equation is solved with the multigrid solvers of \cref{ch:multigrid}
      (\param{gravity_model_fofr_exact_solution}, \param{gravity_model_symmetron_exact_solution},
      \param{gravity_model_dilaton_exact_solution}, \param{gravity_model_dgp_exact_solution},
      \param{gravity_model_cubic_galileon_exact_solution}, \param{gravity_model_horndeski_exact_solution}), as in
      high-resolution codes such as ECOSMOG \cite{Li2012ECOSMOG}.
      For $f(R)$ the field is $u=\sqrt{|f_R|}$ by default, which converges for any $f_{R0}$, $n$ and redshift (the
      original $\ln|f_R|$ can give NaN for small $|f_{R0}|$ at high redshift). The solution of the previous step is the
      starting point of the next solve. The field can be solved on a grid coarser than the force grid (the density
      restricted by a sharp-$k$ cut or CIC, the solution zero above the coarse Nyquist frequency), the uniform-grid
      analogue of Barreira, Bose and Li \cite{BarreiraBoseLi2015}, and only below a chosen redshift (linear or no fifth
      force before). Slow; mostly for testing the approximations.
\end{enumerate}
For the cubic Galileon the Vainshtein equation has no real solution in deep voids ($\delta<-0.6$ today; also found by
Barreira et al.\ \cite{Barreira2013}); the screening factor and the multigrid solver use the zero-discriminant limit
there. With COLA the screened fifth force acts on the residual like the rest of the PM force; the frame uses the linear
(and, for DGP, second-order source) growth.

\paragraph*{Generic Horndeski.} \file{scripts/horndeski_frontend.py} takes $K(\phi,X)$, $G_3(\phi,X)$ and $G_4(\phi)$ as
sympy expressions, derives the background equations from the minisuperspace action, solves them (from today with
$\dd\phi/\dd\ln a$ from the Friedmann constraint, or from early times shooting for $E(1)=1$ for attractors) and writes
$H(a)$, the quasi-static force and the $\alpha$ functions \cite{Bellini2014}. The force (Kimura, Kobayashi and Yamamoto
\cite{Kimura2012}) is the Newtonian force rescaled by $1/(2G_4)$ (running Planck mass) plus a fifth force with the DGP
equation for time-dependent $r_c(a)$ and $\beta(a)$, with Vainshtein screening when $G_{3X}\neq0$, as in Hi-COLA
\cite{Wright2023}; unlike Hi-COLA the 2LPT growth has the theory's second-order source and the initial conditions
are those of the model (\code{horndeski_table} in the solver). K-mouflage is run as in Hi-COLA in the Jordan frame,
with a linear fifth force $G_{\rm eff}/G=A^2(1+2\beta^2/K_X)$ and no kinetic screening.

\begin{validation}[Validation: modified gravity]
\begin{itemize}
\item $\mu$ model: $P_\mu/P_{\rm GR}$ at linear amplitude agrees with $(D_\mu/D_{\rm GR})^2$ from an independent growth
      solution to $0.05\%$ for $k\le0.1\hMpc$ and 0.2--1.2\% up to $0.6\hMpc$ (\file{check_mu_gravity.py}).
\item Symmetron, dilaton, $(m,\beta)$, $f(R)$: the scale-dependent growth of the force agrees with the solver's
      quasi-static growth to $2\times10^{-4}$ (\file{test_quasistatic_models}).
\item Cubic Galileon: linear growth $z=49\to0$ over GR $+2.03\%$ (linear, screening and multigrid alike) against $+2.07\%$
      from the quasi-static equation; at normal amplitude $P/P_{\rm GR}=1.04$--$1.065$ with the exact equation and
      1.03--1.06 with screening, consistent with $\lesssim8\%$ in Barreira et al.
\item Horndeski frontend: reproduces the cubic Galileon ($P(k)$ to $10^{-5}$ in the simulations) and Brans--Dicke's
      $G_{\rm eff}$, $\Sigma$ and $\alpha$'s; K-mouflage background against an independent Einstein-frame integration
      ($2\times10^{-7}$).
\item DGP screening approximation against the multigrid solution: 0.5\% (z = 0) and 1.3\% ($z=1$) rms in $P(k)$ for
      $k<1.5\hMpc$.
\end{itemize}
\end{validation}

\section{Dark-sector interactions}\label{sec:mgsim-darksector}

\paragraph*{Velocity friction (dark scattering).} Elastic scattering off a smooth dark energy (at rest on $N$-body
scales) is a friction of the dark matter, $\dd\bm v/\dd\ln a=-\Gamma\bm v+\bm F$, with
$\Gamma=(1+w)c\,\xi\rho_{\rm DE}/H$ (\code{DarkScattering}) or $\alpha a^3/(\Om_{c0}E)$
(\code{CovariantDarkScattering}) \cite{Simpson2010,BaldiSimpson2015}. It enters the growth equations at all orders,
the COLA frame, the initial velocities and the kick and drift factors of every time-stepping method
(\cref{sec:nbody-timestepping}). Baryons do not scatter: with one CDM+baryon species the fluid feels
$f_c=\Om_{\rm cdm}/\Om_{cb}$ of the drag; with separate baryons the CDM feels the full drag and the baryons none, at the
same steps with their own factors (FastPM from the two-fluid linear growth). Validation (w=$-0.9$, $\xi=50$\,b/GeV,
linear amplitude): the CDM, baryon and CDM+baryon growth suppressions of $-6.16\%$, $-1.31\%$, $-5.42\%$ of the
two-fluid equations are reproduced to $\le0.05\%$ by Quinn, FastPM and COLA. The solver, where the dark energy moves
with the dark matter on large scales, has less suppression for $k\lesssim0.07\hMpc$.

\paragraph*{Coupled quintessence.} In the sub-horizon limit the CDM feels $G(1+2\beta^2)$ from the CDM and $G$ from the
baryons and has the velocity term $+\beta\phi'\bm v$; the baryons feel gravity only. For one CDM+baryon species this is
a momentum-weighted fluid with $G_{\rm eff}/G=1+2\beta^2f_c^2$ and friction $-f_c\beta\phi'$; with separate baryons the
CDM gets its fifth force from the CDM density in a second kick. The CDM density follows the field-dependent mass,
$\rho_ca^3=\Om_{\rm cdm}m(a)/m(1)$. Agreement with the solver's growth: $6\times10^{-4}$ ($\lambda=0.5$, $\beta=0.15$).
The Fading Dark Sector \cite{ObiedVafaWu2026} is \code{cosmology_model = FadingDarkSector} with \code{gravity_model =
FadingDarkSector} (the same code with $\lambda=c$, $\beta=c'$; \code{CoupledQuintessence} also accepts this
cosmology). At the best fit, $c=-0.85$, $c'=0.05$, the CDM feels $G_{\rm eff}/G=1.005$ from the CDM ($1.0036$ for the
one-fluid CDM+baryons today) and the drag $-c'\,\dd\phi/\dd\ln a=0.022$ at $a=1$. Its outputs are bit-identical to
\code{CoupledQuintessence} with the same parameters (regression cases \code{fading_dark_sector_cola},
\code{_fastpm}).

\subparagraph*{The nonlinear boost of the Fading Dark Sector.}\label{sec:mgsims-fds-boost}
\file{SimulationCode/scripts/fds_vs_lcdm_simulations.py} sets up and analyses the simulations that measure the nonlinear
power of the FDS against $\Lambda$CDM and test the halo-model reaction (\cref{sec:hm-react}). The two cosmologies have the
same $A_s$, $n_s$, $h$, $\Omega_b$ and the same matter--radiation equality: the $\Lambda$CDM $\Om_{\rm cdm}$ is the FDS CDM
density at equality scaled as $a^{-3}$ (0.2722 against 0.267 today at the best fit, because the coupled CDM density falls by
2\% after equality), computed from the solver's FDS background. The runs are full N-body (Quinn, no COLA) from the
solver's transfer functions at $z=49$ (2LPT), with fixed amplitudes and paired phases, so the ratio $P_{\rm FDS}/P_{\Lambda\rm
CDM}$ has almost no cosmic variance; two boxes (production: 1000 and $250\Mpch$ with $512^3$ particles and $1024^3$,
$1536^3$ force grids; \code{--short-range} adds P3M in the small box).
\begin{paramfile}[FDS vs LCDM simulations]
python3 scripts/fds_vs_lcdm_simulations.py setup runs/fds   # --preset production|test, --c, --cprime
cd runs/fds && MPIRUN="mpirun -np 64" ./run_all.sh         # solver predictions, then the 8 runs
python3 ../../scripts/fds_vs_lcdm_simulations.py analyse .  # P_FDS/P_LCDM and the boost ratio vs react
\end{paramfile}
\code{run_all.sh} first runs the solver with \code{nonlinear = react} (FDS) and \code{hmcode2020} ($\Lambda$CDM) at the
output redshifts (2, 1, 0.5, 0); \code{analyse} compares $P_{\rm FDS}/P_{\Lambda\rm CDM}$ and
$(P_{\rm NL}/P_L)_{\rm FDS}/(P_{\rm NL}/P_L)_{\Lambda\rm CDM}$ of the simulations with them (up to half the particle Nyquist
frequency, where the power is well above the shot noise). A $64^3$ test (\code{--preset test}) already follows the reaction
to $\sim0.3\%$ in the boost ratio for $k\le0.7\hMpc$ at all four redshifts; the production runs test the scales
$k\approx1$--$3\hMpc$, where the lower concentrations assumed for the FDS halos matter.

\paragraph*{Interacting dark energy.} The exchange is perturbed by the local expansion of the dark matter, so a CDM
particle gains $\frac{\xi}{3}\rho_{\rm DE}$ times the change of its volume and keeps its momentum: particles carry
individual masses updated before each force (the volume from the particle number density, CIC on the particle grid),
and the growth factors include the mass perturbation so the COLA frame follows. This dilution term is most of the
effect (10\% more growth from $z=49$ to 0 for $\xi=-0.1$, $w=-0.9$). Against the solver the extra growth agrees to
$6\times10^{-4}$ for $k<0.02\hMpc$.
